\subsection{可分代数扩张}

定义1. --- 设E是K的一个代数扩张；则称E（在K上）是可分的，如果E的每个在K上有限次的子扩张F都是K上的平展代数（V，第28页，定义1）。

设 E 是 K 的有限次扩张。由于平展代数的每个子代数都是平展的（V，第30页，命题3），假设 E 是 K 的可分扩张，与假设 E 是 K 上的平展代数，这两者是等价的。

命题1. --- 设E是K的一个代数扩张。若E是可分的，则E的每个子扩张E'都是可分的。反之，若E的每个有限次子扩张都是可分的，则E是可分的。

这直接从定义1得出。

命题2. --- 要使域 K 是完全域，必要且充分的是 K 的每个代数扩张都是可分扩张。

首先假设 K 是完全域。由于域是既约环，由引理5（V，第35页）可知，K 的每个有限次扩张都是 K 上的平展代数；因此 K 的每个代数扩张都是可分的。

现在假设 K 是特征为 \(p \neq 0\) 的非完全域。设 \(\Omega\) 是 K 的一个代数闭包。由于 K 是非完全的，存在 \(b \in K\) 不属于 \(K^{p}\) ；令 \(a = b^{1/p}\) 。则扩张 \(\mathrm{K}(a)\) （K 的）是有限次数的 p-根式扩张。由 V，第26页，命题3，存在唯一的将 \(\mathrm{K}(a)\) 映入 \(\Omega\) 的 K-同态，且由于 \([\mathrm{K}(a):\mathrm{K}] > 1\) ，代数 \(\mathrm{K}(a)\) 在 K 上不是平展的（V，第32页，命题4）。换言之，K 的有限次扩张 \(\mathrm{K}(a)\) 不是可分的。

推论。--- 特征为0的域或有限域的每一个代数扩张都是可分扩张。

这由V，第7页，命题5得出。

